\chapter*{Abstract}
\addcontentsline{toc}{chapter}{Abstract}

This thesis presents the \textbf{hardware side} of an integrated effort
to deploy and document community ICT infrastructure for
low-connectivity, resource-limited contexts. Conducted jointly with a
companion thesis on the software and services side {[}Motje, 2026{]},
the work delivers three contributions: a reproducible community wireless
network recipe based on OpenWrt commodity access points with 802.11s
mesh backhaul; a refurbished-endpoint pipeline that provisions donated
business laptops through PXE-booted Clonezilla and a curated Linux Mint
golden-master image; and an open, web-and-PDF, GitHub-hosted
\textbf{Community-Network-Handbook} consolidating the operational
knowledge into a living artefact. The three contributions were validated
through a field deployment at N Mutschuana Primary School, Gochas,
Namibia, in March 2026: a seven-AP mesh covering the school and nine
refurbished Lenovo ThinkPads provisioned in approximately one hour.
Validation is reported through three qualitative lenses ---
\emph{coverage}, \emph{sufficiency}, \emph{adaptation} --- and the
eleven operational lessons produced were committed back into the
handbook during the deployment itself.

\textbf{Keywords:} community networks, ICT for development, refurbished
hardware, OpenWrt, mesh networking, PXE, Clonezilla, knowledge
management, AUCOOP, Namibia, open documentation, Diátaxis.

\hypertarget{notes}{%
\subsection{Notes}\label{notes}}

I dont like the approach of the abstract, to be reviewed when the
document is completed. the way is introduced in Jaume's thesis is a
better approach
